\section{Conclusion}
\label{sec:main-conclusion}

We formulated concept updating as a risk-sensitive test-time decision: a representation should be refined only when the induced information split has positive predictive value after estimation error and complexity costs are paid. This yields two structural consequences: a repair-only memory cannot create distinctions absent from its observable state, and a valid refinement must be guarded by an empirical certificate tied to the downstream loss. Conceptual Attention instantiates the principle in the main experiments as a fixed-bank sequence layer that reads from learned concept identities and updates their value memory during the forward pass; the implementation evaluates the exact streaming recurrence by chunked affine scan. The predictive CoM prototype sharpens the test-time-learning claim: its fast state is a dense, gradient-free sufficient-statistic repair of observed consequences. The current training probes suggest that the main practical bottleneck is not sparse routing or a write gate, but the consequence-feature capacity and decoder used by this state. The appendix gives the graph-growing refinement extension. The remaining empirical question is quantitative: after the final zero-shot and same-corpus PPL evaluations are filled, CoM either matches the relevant recurrent baselines under the same protocol or exposes the concrete optimization and systems gaps that the theory does not by itself remove.
